\subsection{Gleichwertige Geradendarstellungen}
\label{subsec:analytic_geometry_equivalent_line_representations}

Eine Gerade besitzt keine eindeutige Parameterdarstellung. Das Erkennen
gleichwertiger Darstellungen verhindert, dass verschiedene Formeln
fälschlich als verschiedene Geraden interpretiert werden.

\begin{KeyPointBox}{Dieselbe Gerade unterschiedlich darstellen}{analytic_geometry_equivalent_lines}
  Dieselbe Gerade entsteht, wenn
  \begin{itemize}[leftmargin=1.6em,itemsep=0.4em]
    \item ein anderer Punkt der Geraden als Aufpunkt gewählt wird und
    \item der Richtungsvektor durch ein beliebiges von null verschiedenes
      Vielfaches ersetzt wird.
  \end{itemize}
\end{KeyPointBox}

\noindent
\begin{minipage}[t]{0.32\linewidth}
  \centering
  \begingroup
    \def\LineUnit{0.62}
    \def\LineSupportX{2}\def\LineSupportY{2}
    \def\LineDirectionX{1}\def\LineDirectionY{-1}
    \def\LineLambdaMin{-2}\def\LineLambdaMax{2}
    \CachedTikzFigure{content/resources/chapters/analytic_geometry/tikz/parameterized_line_representation.tex}
  \endgroup
\end{minipage}\hfill
\begin{minipage}[t]{0.32\linewidth}
  \centering
  \begingroup
    \def\LineUnit{0.62}
    \def\LineSupportX{1}\def\LineSupportY{3}
    \def\LineDirectionX{1}\def\LineDirectionY{-1}
    \def\LineLambdaMin{-1}\def\LineLambdaMax{3}
    \CachedTikzFigure{content/resources/chapters/analytic_geometry/tikz/parameterized_line_representation.tex}
  \endgroup
\end{minipage}\hfill
\begin{minipage}[t]{0.32\linewidth}
  \centering
  \begingroup
    \def\LineUnit{0.62}
    \def\LineSupportX{2}\def\LineSupportY{2}
    \def\LineDirectionX{-1}\def\LineDirectionY{1}
    \def\LineLambdaMin{-2}\def\LineLambdaMax{2}
    \CachedTikzFigure{content/resources/chapters/analytic_geometry/tikz/parameterized_line_representation.tex}
  \endgroup
\end{minipage}

\[
  \begin{aligned}
    g_1:\quad \vec x&=\begin{pmatrix}2\\2\end{pmatrix}
      +\lambda\begin{pmatrix}1\\-1\end{pmatrix},\\[0.7em]
    g_2:\quad \vec x&=\begin{pmatrix}1\\3\end{pmatrix}
      +\mu\begin{pmatrix}1\\-1\end{pmatrix},\\[0.7em]
    g_3:\quad \vec x&=\begin{pmatrix}2\\2\end{pmatrix}
      +\nu\begin{pmatrix}-1\\1\end{pmatrix}.
  \end{aligned}
\]

\begin{ExampleBox}{Geradendarstellungen vergleichen}{analytic_geometry_compare_line_representations}
  Der Aufpunkt von $g_2$ liegt auf $g_1$, denn für $\lambda=-1$ gilt
  \[
    \begin{pmatrix}2\\2\end{pmatrix}
    -\begin{pmatrix}1\\-1\end{pmatrix}
    =\begin{pmatrix}1\\3\end{pmatrix}.
  \]
  Zudem ist der Richtungsvektor von $g_3$ das $(-1)$-Fache des
  Richtungsvektors von $g_1$. Somit beschreibt jede der drei Gleichungen
  \Cloze{5cm}{dieselbe Gerade}.
\end{ExampleBox}
